\documentclass[12pt,letterpaper, onecolumn]{exam}
\usepackage{amsmath}
\usepackage{amssymb}
\usepackage{amsthm}
\usepackage{graphicx}
\usepackage{mathrsfs}
\usepackage{xcolor}
\usepackage[lmargin=71pt, tmargin=1.2in]{geometry}  %For centering solution box
\usepackage[english]{babel}
\newtheorem{theorem}{Theorem}
\newtheorem{lemma}[theorem]{Lemma}
\usepackage{accents}
\newcommand{\ubar}[1]{\underaccent{\bar}{#1}}
\newcommand{\bs}[1]{\boldsymbol{#1}}
\newcommand{\bm}[1]{\mathbf{#1}}
\newcommand{\der}[2]{\frac{d #1}{d#2}}
\newcommand{\pder}[3]{\frac{\delta^{#3} #1}{\delta^{#3}#2}}
\newcommand{\rs}{$X_1, ..., X_n$}
\newcommand{\nsum}{\sum_{i=1}^n}
\newcommand{\nprod}{\prod_{i=1}^n}
\newcommand*{\Comb}[2]{{}^{#1}C_{#2}}%
\newcommand{\ifc}{\textrm{ if }}
\newcommand{\real}{\mathbb{R}}
\newcommand{\nat}{\mathbb{N}}
\newcommand{\power}{\mathcal{P}}
\newcommand{\rational}{\mathbb{Q}}
\newcommand{\cutelim}{\lim\limits_{n\to\infty}}
\newcommand{\salg}{$\sigma$-algebra }
\newcommand{\calA}{\mathcal{A}}
\newcommand{\calO}{\mathcal{O}}
\newcommand{\calC}{\mathcal{C}}
\newcommand{\calB}{\mathcal{B}}
\newcommand{\calI}{\mathcal{I}}
\newcommand{\calF}{\mathcal{F}}
\newcommand{\calG}{\mathcal{G}}
\newcommand{\gensig}[1]{\sigma\langle #1 \rangle} % notation for a sigma algebra generated by some set 
\newcommand{\om}{\mu^*} % outer measure 
\newcommand{\rat}{\mathbb{Q}}
\newcommand{\tribrac}[1]{\langle #1 \rangle}
\newcommand{\borel}{\mathcal{B}(\mathbb{R})}
\newcommand{\boreltwo}{\mathcal{B}(\mathbb{R}^2)}
\newcommand{\borelk}[1]{\mathcal{B}(\mathbb{R}^#1)}
\newcommand{\probspace}{(\Omega, \calF, P)}

\newcommand{\Iunion}{\bigcup_{\alpha \in I}}
\newcommand{\Iintersect}{\bigcap_{\alpha \in I}}

% \chead{\hline} % Un-comment to draw line below header
\thispagestyle{empty}   %For removing header/footer from page 1

\begin{document}

\begingroup  
    \centering
    \LARGE STAT 6410\\
    \LARGE Homework \\[0.5em]
    \large Sarah Kilgard\par
\endgroup
\rule{\textwidth}{0.4pt}
\pointsdroppedatright   %Self-explanatory
\printanswers
\renewcommand{\solutiontitle}{\noindent\textbf{Ans:}\enspace}   %Replace "Ans:" with starting keyword in solution box

\begin{questions}
%%%%%%%%%%% Question 1 %%%%%%%%%%%%%%%%%%%%    
    \question Chapter 2, Problem 2.1

    Let $\Omega_i$, $i = 1, 2$ be two nonempty sets and $T: \Omega_1 \to \Omega_2$ be a map. Then for any collection $\{A_\alpha: \alpha \in I\}$ of subsets of $\Omega_2$, show that: 
    \[T^{-1}\left(\bigcup_{\alpha \in I} A_\alpha\right) = \bigcup_{\alpha \in I} T^{-1}(A_\alpha)\]
    \[\text{and } T^{-1}\left(\bigcap_{\alpha \in I}A_\alpha\right) = \bigcap_{\alpha \in I} T^{-1}(A_\alpha)\]
    Further, $(T^{-1}(A))^c = T^{-1}(A^c)$ for all $A \subset \Omega_2$. (These are known using \textit{de Morgan's} laws.)
    \begin{solution}
    \begin{proof}
        \textbf{Part 1:} WTS: $T^{-1}\left(\bigcup_{\alpha \in I} A_\alpha\right) = \bigcup_{\alpha \in I} T^{-1}(A_\alpha)$

        \underline{Part 1.a:} $T^{-1}\left(\bigcup_{\alpha \in I} A_\alpha\right) \subset \bigcup_{\alpha \in I} T^{-1}(A_\alpha)$ \\
        Let $\omega \in T^{-1}\left(\bigcup_{\alpha \in I} A_\alpha\right)$ be given. \\
        Thus, $T(\omega) \in \Iunion A_\alpha \: \implies \: \exists \beta \in I \: s.t. \: T(\omega) \in A_\beta \text{ and } \omega \in T^{-1}(A_\beta).$\\
        $T^{-1}(A_\beta) \subset \Iunion T^{-1}(A_\alpha) \: \implies \omega \in \Iunion T^{-1}(A_\beta)$

        Thus, $T^{-1}\left(\bigcup_{\alpha \in I} A_\alpha\right) \subset \bigcup_{\alpha \in I} T^{-1}(A_\alpha)$
        
        \underline{Part 1.b:} $T^{-1}\left(\bigcup_{\alpha \in I} A_\alpha\right) \supset \bigcup_{\alpha \in I} T^{-1}(A_\alpha)$

        Let $\omega \in \Iunion T^{-1}(A_\alpha)$ be given.\\
        Thus, $\exists \: \beta \in I \: s.t. \: T(\omega) \in A_\beta $ and $\omega \in T^{-1}(A_\beta)$. \\
        Since, $A_\beta \subset \Iunion A_\alpha, \: \omega \in T^{-1}(\Iunion A_\alpha)$

        Thus, $T^{-1}\left(\bigcup_{\alpha \in I} A_\alpha\right) = \bigcup_{\alpha \in I} T^{-1}(A_\alpha)$.

        \textbf{Part 2:} WTS: $T^{-1}\left(\bigcap_{\alpha \in I} A_\alpha\right) = \bigcap_{\alpha \in I} T^{-1}(A_\alpha)$

        \underline{Part 2.a:} $T^{-1}\left(\bigcap_{\alpha \in I} A_\alpha\right) \subset \bigcap_{\alpha \in I} T^{-1}(A_\alpha)$ \\
        Let $\omega \in T^{-1}(\Iintersect A_\alpha)$ be given. Thus, $T(\omega) \in \Iintersect A_\alpha$. \\
        $\implies \: \forall \alpha \in I, \: T(\omega) \in A_\alpha$. $\therefore \omega \in T^{-1}(A_\alpha) \: \forall \alpha \in I \: \implies \: \omega \in \Iintersect T^{-1}(A_\alpha)$
        
        \underline{Part 2.b:} $T^{-1}\left(\bigcap_{\alpha \in I} A_\alpha\right) \supset \bigcap_{\alpha \in I} T^{-1}(A_\alpha)$ \\
        Let $\omega \in \Iintersect T^{-1}(A_\alpha)$ be given. Thus, $\forall \: \alpha \in I, \: \omega \in T^{-1}(A_\alpha)$ and $T(\omega) \in A_\alpha$. \\
        $T(\omega) \in \Iintersect A_\alpha \: \implies \: \omega \in T^{-1}(\Iintersect A_\alpha)$

        Thus, $T^{-1}\left(\bigcap_{\alpha \in I} A_\alpha\right) = \bigcap_{\alpha \in I} T^{-1}(A_\alpha)$.

        \textbf{Part 3:} WTS: $(T^{-1}(A))^c = T^{-1}(A^c)$ for all $A \subset \Omega_2$

        Let $A \subset \Omega_2$ be given. 

        \underline{Part 3.a:} WTS $(T^{-1}(A))^c \subset T^{-1}(A^c)$ \\
        Let $\omega \in (T^{-1}(A))^c$ be given. \\
        Thus, $\omega \notin T^{-1}(A) \: \implies \: T(\omega) \notin A \: \implies T(\omega) \in A^c$\\
        $\therefore \: \omega \in T^{-1}(A^c)$

        \underline{Part 3.b:} WTS $(T^{-1}(A))^c \supset T^{-1}(A^c)$ \\
        Let $\omega \in T^{-1}(A^c)$ be given. \\ Thus, $T(\omega) \in A^c \: \implies T(\omega) \notin A$. \\
        $\implies \: \omega \notin T^{-1}(A) \: \implies \: \omega \in (T^{-1}(A))^c$

        Thus, $(T^{-1}(A))^c = T^{-1}(A^c)$ for all $A \subset \Omega_2$
    \end{proof}
    \end{solution}

%%%%%%%%%%% Question 2 %%%%%%%%%%%%%%%%%%%%    
    \question Chapter 2, Problem 2.3

    Let $f, g: \Omega \to \real$ be $\tribrac{\calF, \borel}$-measurable. Set
    \[h(\omega) = \frac{f(\omega)}{g(\omega)}I(g(\omega) \neq 0), \: \omega \in \Omega.\]
    Verify that $h$ is $\tribrac{\calF, \borel}$-measurable. 
    \begin{solution}
    \begin{proof}
        $\ubar{x} \equiv (x_1, x_2) \in \real^2$. \\
        Define: $h_1(\ubar{x}) \equiv \frac{x_1}{x_2}I(x_2 \neq 0)$. \\
        $h_{1}^{-1}((-\infty, a]) \equiv \{(x_1, x_2) \in \real^2: x_2 = 0, \: a \leq 0\} \\ \cup \{(x_1, x_2) \in \real^2: x_2 > 0, \: x_1 - a x_2  \leq 0\} \cup \{(x_1, x_2) \in \real^2: x_2 < 0, \: x_1 - a x_2 \geq 0\}$
        
        $h_{1}^{-1}((-\infty, a]) \in \borelk{2}$ for all $a \in \real$. Thus, $h_1$ is $\tribrac{\borelk{2}, \borel}-$measurable, by Definition 2.1.2.

        By Proposition 2.1.3, $\eta = (f, g)$ is $\tribrac{\calF, \borelk{2}}$-measurable.

        By Proposition 2.1.1, $h = \nu \circ h_1$ is $\tribrac{\calF, \borel}$-measurable.
    \end{proof}
    \end{solution}

%%%%%%%%%%% Question 3 %%%%%%%%%%%%%%%%%%%%    
    \question Chapter 2, Problem 2.6

    Let $X_i$, $i = 1, 2, 3$ be random variables on a probability space $(\Omega, \calF, P)$. Consider the random equation (in $t \in \real$):
    \[X_1(\omega)t^2 + X_2(\omega)t + X_3(\omega) = 0. \quad \quad (6.1)\]
    \begin{itemize}
        \item[(a)] Show that $A \equiv \{\omega \in \Omega: \text{Equation (6.1) has two distinct real roots}\} \in \calF.$
        \item[(b)] Let $T_1(\omega)$ and $T_2(\omega)$ denote the two roots of (6.1) on A. Let 
        \[f_i(\omega) = \begin{cases}
            T_i(\omega) & \text{on } A \\
            0 & \text{on } A^c
        \end{cases},\]
        $i = 1, 2$. Show that $(f_1, f_2)$ is $\tribrac{\calF, \boreltwo}$-measurable.
    \end{itemize}
    \begin{solution}
        \begin{itemize}
            \item[(a)] 
            \begin{proof}
            We can rewrite (without changing the set composition) A as: \\ $A = \{\omega \in \Omega: X_2(\omega)^2 - 4 X_1(\omega) X_3(\omega) > 0\} \cap \{\omega \in \Omega: X_1(\omega) \neq 0\}$. This pulls from the definition of the determinant of a quadratic equation from elementary algebra.

            Because $X_1, X_2, X_3$ are random variables, they are $\tribrac{\calF, \borel}$- measurable transformations from $\Omega$ to $\real$. \\
            By Corollary 2.1.4., since $X_1, X_2, X_3$ are $\tribrac{\calF, \borel}$- measurable functions from $\Omega$ to $\real$, this collection is closed under addition, multiplication, and scalar multiplication. Thus, $X_2^2 - 4X_1 X_3$ is $\tribrac{\calF, \borel}$- measurable (and a random variable on the probability space $(\Omega, \calF, P)$). 

            Thus, by the definition of a random variable, $A \in \calF$.
            \end{proof}
            \item[(b)] 
            \begin{proof}
            Again, referring to elementary algebra, we can define $T_1$ and $T_2$ as:
            \[T_1(\omega) = \frac{-X_2(\omega) - \sqrt{X_2^2(\omega) - 4 X_1(\omega) X_3(\omega)}}{2X_1(\omega)}\] 
            \[ T_2(\omega) = \frac{-X_2(\omega) + \sqrt{X_2^2(\omega) - 4 X_1(\omega) X_3(\omega)}}{2X_1(\omega)} \]

            Define $f_i$:
            \[f_i(\omega) = T_i(\omega) I(\omega \in A) = T_i(\omega)I(X_1(\omega) \neq 0) I(X_2^2(\omega) - 4 X_1(\omega) X_3(\omega) > 0)\]

           $g(\omega) \equiv \sqrt{X_2^2(\omega) - 4 X_1(\omega) X_3(\omega)} I(X_2^2(\omega) - 4 X_1(\omega) X_3(\omega) > 0)$

           $g: \Omega \to \real$ is piecewise continuous on $\Omega$. Thus, by problem 2.2, $g$ is $\tribrac{\calF, \borel}$-measurable. 

           $h(\omega) = -X_2(\omega)$ and $l(\omega) = 2X_1(\omega)$ are both $\tribrac{\calF, \borel}$-measurable by Corollary 2.1.4.

           By, Corollary 2.1.4, $h(\omega) - g(\omega)$ and $h(\omega) + g(\omega)$ are both $\tribrac{\calF, \borel}$-measurable.

           Finally, by Problem 2.3, \\$f_1(\omega) = \frac{h(\omega) - g(\omega)}{l(\omega)}I(l(\omega) \neq 0) = T_2(\omega) I(\omega \in A)$ and \\
           $f_2(\omega) = \frac{h(\omega) + g(\omega)}{l(\omega)}I(l(\omega) \neq 0) = T_1(\omega) I(\omega \in A)$ \\
           are $\tribrac{\calF, \borel}$-measurable.
            \end{proof}
        \end{itemize}
    \end{solution}

%%%%%%%%%%% Question 4 %%%%%%%%%%%%%%%%%%%%    
    \question Extra Problem A

    Let $X: \Omega_1 \to \Omega_2$ be a $\tribrac{\calF_1, \calF_2}$ be a measurable transformation. Consider the collection (of subsets of $\Omega_1$) defined by $\gensig{X} = \{X^{-1}(B), B \in \calF_2\}$. Then show that
    \begin{itemize}
        \item[(a)] $\gensig{X}$ is \salg,
        \item[(b)] $\gensig{X}$ is the \textit{smallest} \salg that makes $X$ measurable, (i.e. if X is $\tribrac{\calG, \calF_2}$-measurable, then $\gensig{X} \subset \calG)$.
    \end{itemize}
    \begin{solution}
        \begin{itemize}
            \item[(a)] 
            \begin{proof}
                \textbf{Part 1:} WTS: $\Omega_1 \in \gensig{X}$. \\
                This is trivial. 

                \textbf{Part 2:} WTS: $A \in \gensig{X} \: \implies \: A^c 
                \in \gensig{X}$

                Let $A \in \gensig{X}$ be given. \\ 
                Thus, there exists some $B \in \calF_2$ such that $A = X^{-1}(B)$. \\
                $A^c = (X^{-1}(B))^c = X^{-1}(B^c)$, by Problem 2.1. \\
                $B^c \in \calF_2$ since $\calF_2$ is a \salg. \\
                Thus, $A = X^{-1}(B^c) \in \gensig{X}$

                \textbf{Part 3:} WTS $A_1, A_2, ..., A_n \in \gensig{X} \: \implies \Iunion A_\alpha \in \gensig{X}$, where $I = (1, 2, ... , n)$ for some fixed n

                Let $A_1, A_2, ..., A_n \in \gensig{X}$ be given. Thus there exists $B_i \in \calF_2$ such that $A_\alpha = X^{-1}(B_\alpha)$ for all $\alpha \in I$. \\
                $\Iunion A_\alpha = \Iunion X^{-1}(B_\alpha) = X^{-1}\left(\Iunion B_\alpha\right)$, by Problem 2.1\\
                $\Iunion B_\alpha \in \calF_2$ since $\calF_2$ is a \salg. \\
                Thus, $\Iunion A_\alpha = X^{-1}\left(\Iunion B_\alpha\right) \in \gensig{X}$. 

                Thus, $\gensig{X}$ is a \salg.
            \end{proof}
            \item[(b)] 
            \begin{proof}
                Suppose $X$ is $\tribrac{\calG, \calF_2} $-measurable. \\
                Thus, for all $B \in \calF_2$, $X^{-1}(B) \in \calG \: \implies \: \gensig{X} = \{X^{-1}(B), B \in \calF_2\} \subset \calG$.

                Thus, $\gensig{X}$ is the \textit{smallest} \salg that makes $X$ measurable.
            \end{proof}
        \end{itemize}
    \end{solution}

%%%%%%%%%%% Question 5 %%%%%%%%%%%%%%%%%%%%    
    \question Extra Problem B

    Give an example of a function $f$ that is \textbf{not} $\tribrac{\borel, \borel}$-measurable.
    \begin{solution}
    From the results of Lebesgue Measure on $\real$, we know that there exists $A \in \real$ such that $B \notin \borel$.

    $f: (\real, \borel) \to (\real, \borel)$ 

    As such, 
    $f(\omega) = \begin{cases}
        1 & \omega \in B \\
        0 & \omega \notin B
    \end{cases}$  
    is not $\tribrac{\borel, \borel}$-measurable.

    Indicator functions are only measurable if $A$ is in $\borel.$
    \end{solution}
    
\end{questions}
\end{document}